% Thesis English abstract — current verified state as of 18 November 2026.
\begin{kbtuabstract}{ABSTRACT}
This diploma project addresses the development of a control system for an autonomous mobile
manipulator intended to assist agricultural operations in hydroponic greenhouses. The system
combines an omnidirectional mobile platform, a Schneider Electric Lexium collaborative robot,
a four-finger compliant gripper, an RGB-D camera and modular robot-control software. The
proposed operating pipeline links object perception and three-dimensional localization with
safe approach, grasping, sorting and placement. Current work covers robot modeling, a
physics-based digital twin, perception experiments, trajectory checks and preliminary
hardware integration. In NVIDIA Isaac Sim, a hardware-constrained validation suite passed
five of five scenarios, including startup, single-fruit handling, three-fruit sorting,
randomized reachable targets and safe rejection/recovery. The simulated single-fruit test
reported a localization error of 1.97~mm and a placement error of 12.81~mm; a separate
three-fruit scenario successfully handled three of three targets. A further five-object
sorting demonstration completed five of five placements in its test scene. These results
are simulation-level evidence; they do not constitute validation of autonomous greenhouse
harvesting on the physical robot. The next stages are end-effector and camera calibration,
physical grasp-force testing, navigation integration and repeatable greenhouse trials.
\end{kbtuabstract}
